% Transcription — fonds Alexandre Grothendieck, shelfmark 49, batch 1 (pages 1-9).
% Established from the facsimile archives/batches/49/batch-01.pdf (a mirror of
% https://grothendieck.umontpellier.fr/49.pdf), per the skill
% transcribe-grothendieck. The folder is nine pages and this is its only batch.
% It is one continuous run in one hand: page 1 is a cover bearing only the title
% "Picard des schémas abéliens", so it carries no \page{} and its words open the
% file as the section heading; pages 2 to 9 are the run itself, and every one of
% them is mathematics. No unrelated verso anywhere in the folder, which is
% unusual for this fonds.
% The run is two theorems with their corollaries. The first (p. 2) is that
% Pic^0 of a projective abelian scheme is again an abelian scheme, proved
% through phi_L and the theorem of Hopf; its Corollaire "Main 1" (p. 4) is that
% R^*f_*(O_A) is the exterior algebra on R^1f_*(O_A). The second (p. 6) reads
% the arrow Pic_{A/S} -> Hom_{S-gr}(A,A') and identifies its kernel as the
% sections satisfying pi^*(xi) = pr_1^*(xi) + pr_2^*(xi); pages 8 and 9 are its
% proof, which base-changes to S_1 = A.
% The hand is fast throughout. Following folder 29 batches 8 and 9, the
% mathematical statements are transcribed in full and the connective prose is
% left \ill{} wherever the leaf will not support a reading — the sentences
% between the displays are the palimpsest part of these pages, and several of
% them, including most of the diagonal note up the left margin of page 2, are
% not recoverable at the resolution the film carries.
% Pass: Opus 5 (claude-opus-5), 2026-08-30 — first pass, unchecked against the pages by a human.
\documentclass[11pt,a4paper]{article}
% Le préambule commun à toutes les transcriptions du fonds Grothendieck.
%
% Il tient en une idée : **l'incertitude se compose, elle ne se résout pas.**
% Une transcription qui lisse un mot douteux a détruit la seule chose pour
% laquelle on la faisait — un lecteur doit pouvoir savoir, à chaque endroit,
% ce qui était sur le feuillet et ce qui a été ajouté. Les six macros qui
% suivent sont donc l'appareil critique, et non de la décoration.
%
% Les mêmes commandes sont reconnues par `scripts/render.mjs`, qui produit la
% vue de lecture du site. Une macro ajoutée ici doit l'être aussi là-bas,
% faute de quoi le rendu échouera bruyamment — ce qui est voulu : un rendu
% silencieusement faux est pire qu'un rendu absent.

\ProvidesPackage{grothendieck}[2026/08/08 Transcriptions du fonds Grothendieck]

\usepackage{amsmath,amssymb,amsthm}
\usepackage{mathrsfs}
\usepackage{stmaryrd}
% Les diagrammes commutatifs sont partout dans ce fonds : c'est la langue
% même de la géométrie algébrique de Grothendieck, et une transcription qui
% les rendrait en prose ne transcrirait rien.
\usepackage{tikz-cd}
% Français par défaut — c'est la langue des feuillets — mais l'anglais est
% chargé aussi : la traduction et le résumé sont des documents anglais, et la
% typographie française (espaces avant les deux-points) y serait une faute.
% Ces documents basculent par \selectlanguage{english} en tête de corps.
\usepackage[english,french]{babel}
\usepackage{fontspec}
\usepackage{geometry}
\geometry{a4paper,margin=25mm,marginparwidth=28mm}
\usepackage{marginnote}
\usepackage{xcolor}
% ulem, not soul. `soul`'s \st and \ul are built on a character-by-character
% scan that XeLaTeX's font machinery breaks: the document fails with an
% undefined control sequence rather than degrading. ulem does the same two jobs
% and is Unicode-engine safe. `normalem` stops it hijacking \emph, which the
% transcriptions use for Grothendieck's own emphasis.
\usepackage[normalem]{ulem}
\usepackage{hyperref}
\hypersetup{colorlinks=true,linkcolor=black,urlcolor=black,pdfborder={0 0 0}}

% --- Métadonnées du lot -----------------------------------------------------
% Reprises telles quelles par le script de rendu, qui les affiche en tête de
% la vue de lecture. Elles fixent ce que le fichier prétend couvrir : sans
% elles, une transcription flotte sans point d'attache dans un fonds de seize
% mille feuillets.

\newcommand{\folder}[1]{\gdef\grothFolder{#1}}
\newcommand{\batch}[1]{\gdef\grothBatch{#1}}
\newcommand{\leaves}[2]{\gdef\grothFirst{#1}\gdef\grothLast{#2}}
\newcommand{\foldertitle}[1]{\gdef\grothFoldertitle{#1}}
\gdef\grothFolder{?}\gdef\grothBatch{?}\gdef\grothFirst{?}\gdef\grothLast{?}\gdef\grothFoldertitle{}

% --- Le résumé --------------------------------------------------------------
%
% Il ouvre la lecture modernisée, sous le titre « Résumé », et il est composé
% en petit. Ce n'est pas un ornement : c'est le seul passage du document qui
% ne soit pas des mathématiques, et un lecteur doit pouvoir le reconnaître
% avant de l'avoir lu — puis le sauter s'il connaît déjà le sujet, ou s'y
% arrêter s'il ne le connaît pas. Le filet qui le suit marque la reprise du
% corps.

\newenvironment{resume}
  {\small\setlength{\parskip}{0.45em}}
  {\par\medskip\hrule\medskip}

% --- Mots-clés ---------------------------------------------------------------
%
% En anglais, en clôture du résumé de la lecture modernisée : le vocabulaire
% moderne sous lequel chercher ce que les feuillets construisent. Le manifeste
% du site (`npm run manifest`) les extrait pour étiqueter la cote entière —
% cette ligne est la source unique des tags, il n'y a pas de fichier à côté.
\newcommand{\keywords}[1]{\par\smallskip\noindent
  {\footnotesize\textsc{Keywords} --- \textit{#1}}\par}

% --- L'appareil critique ----------------------------------------------------

% Le numéro de feuillet, en marge. C'est l'ancre qui permet de régler un
% désaccord sur une formule en regardant *un* feuillet plutôt qu'une chemise
% de six cents. Le site s'en sert aussi pour tourner les pages du fac-similé
% quand on fait défiler la transcription.
\newcommand{\leaf}[1]{%
  \marginnote{\footnotesize\color{gray!70}\textsf{#1}}%
  \label{leaf:#1}\ignorespaces}

% Illisible : on ne devine pas. Un mot inventé qui se lit comme les autres est
% le pire résultat possible.
\newcommand{\ill}{\textcolor{red!60!black}{[\,\ldots\,]}}

% Lecture douteuse : le mot est donné, et signalé comme incertain.
\newcommand{\uncertain}[1]{\uline{#1}}

% Ajout de l'éditeur — une lettre manquante, un mot sous-entendu.
\newcommand{\add}[1]{\textcolor{blue!50!black}{[#1]}}

% Biffé par Grothendieck lui-même. Ce qu'il a rayé fait partie du document :
% une rature dit souvent où la pensée a changé de direction.
\newcommand{\struck}[1]{\sout{#1}}

% Note du transcripteur, sur l'état matériel du feuillet.
\newcommand{\note}[1]{\par\smallskip{\small\color{gray!60!black}\textit{[#1]}}\par\smallskip}

% Note marginale de Grothendieck, distincte du corps du texte.
\newcommand{\marginal}[1]{\par\smallskip\noindent\hspace{1em}%
  {\small\color{gray!60!black}#1}\par\smallskip}

% --- Titre ------------------------------------------------------------------

\AtBeginDocument{%
  \begin{center}
    {\large\bfseries Fonds Alexandre Grothendieck --- cote n\textsuperscript{o}~\grothFolder}\\[2pt]
    {\itshape\grothFoldertitle}\\[4pt]
    {\small feuillets \grothFirst--\grothLast\ (lot \grothBatch)}\\[2pt]
    {\footnotesize\color{gray!60!black}Transcription établie d'après le fac-similé numérisé
     par l'Université de Montpellier.\\
     Les lectures incertaines sont soulignées, les illisibles notées [\,\ldots\,],
     les ajouts entre crochets.}
  \end{center}
  \vspace{1em}}

\endinput


\folder{49}
\batch{1}
\pages{1}{9}
\dating{s.d.}
\watermark{Édition de démonstration}
\foldertitle{Picard des schémas abéliens : notes manuscrites (s.d.)}

\begin{document}

\section*{Picard des schémas abéliens}

\note{titre de sa main, souligné, seul objet du premier feuillet, qui ne porte
rien d'autre. Le feuillet est une chemise et ne reçoit donc pas de \texttt{page},
conformément à l'usage suivi depuis le dossier 29 : le saut dans la numérotation
est la trace qu'il a été passé}

\page{2}

\note{sur toute la suite il écrit $A'$ pour le schéma abélien dual et $A_{0}$,
$A'_{0}$ pour les fibres en un point. Le foncteur de Picard est tantôt indicé
$A/S$, tantôt $X/S$, sans qu'aucune distinction ne paraisse voulue : les deux
lettres désignent le même schéma abélien et le $X$ est laissé tel quel là où il
paraît}

\textbf{Théorème}\quad Soit $A$ un schéma abélien projectif sur $S$. Alors
\struck{$\mathrm{Pic}^{\tau}_{A/S}$} $\mathrm{Pic}^{0}_{A/S}$ \struck{est} un
schéma abélien sur $S$, ouvert et fermé dans $\mathrm{Pic}_{A/S}$, ayant même
dimension relative sur $S$ que $A$.

Comme $A$ est simple sur $S$ et que les fibres ont des Nér.-Sév. sans torsion,
on sait que $\mathrm{Pic}^{\tau}_{A/S}$ est propre sur $S$, et fermé dans
$\mathrm{Pic}_{A/S}$. Il est en tous cas \emph{ouvert} dans
$\mathrm{Pic}^{\tau}_{A/S}$. Il reste à montrer qu'il est \emph{simple} sur $S$
(on sait déjà que ses fibres sont universellement connexes), et
\struck{la bonne dimension}.

Soit donc $L$ un Module ample sur $A$ rel. $S$, il définit une section $s$ de
$\mathrm{Pic}_{A/S}$ sur $S$, et un morphisme

\struck{\ill{}}
\struck{$P = s + \mathrm{Pic}^{0}_{A/S} \subset \mathrm{Pic}_{A/S}$. Et \ill{} un
morphisme}

\note{deux lignes biffées d'un trait continu ; la seconde se laisse lire en
partie, la première pas}

\[
  \varphi_{s} : A \longrightarrow A'\struck{^{0}} = \mathrm{Pic}^{\tau}_{X/S}
\]

\note{l'exposant $0$ de $A'$ est biffé, et l'indice du foncteur de Picard est
bien $X/S$ et non $A/S$ ; voir la note d'ouverture de la page}

obtenu posant
\[
  \varphi_{s}(t) = \text{classe de } \tau_{t}^{*}(L)\, L^{-1}
\]

\note{le $t$ en indice de $\tau$ est souligné, comme il le fait pour une section
d'un $T$-point}

[Comme $\varphi_{s}$ transforme section unité en section unité, c'est un hom.\
de schémas en groupes (cf.).] \struck{Donc} L'hom.\ $A_{0} \to A'_{0}$ défini
par $\varphi$ est \textbf{surjectif} [car le noyau est fini d'après le critère
d'amplitude de Weil, d'autre part
\[
  \dim A'_{0} \;\leqslant\; \dim H^{1}(A_{0}, \mathcal{O}_{A_{0}})
  \;\leqslant\; \dim A_{0}
\]

\note{sous la première majoration il écrit une accolade et, dessous,
\ill{} \guillemotleft{} de Pic \guillemotright{} — les deux ou trois mots qui
précèdent \emph{de Pic} ne se laissent pas lire ; ce qu'ils justifient est que
l'espace tangent de $\mathrm{Pic}$ à l'origine est $H^{1}(A_{0},
\mathcal{O}_{A_{0}})$. Sous la seconde, deux lignes qui se chevauchent :
\guillemotleft{} th.\ de Hopf $+$ \uncertain{nullité} de $H^{i}(A_{0},
\mathcal{O}_{A_{0}})$ p\add{ou}r $i > \dim A_{0}$ \guillemotright{} — soit
l'argument classique : $H^{*}$ est une algèbre de Hopf, donc l'algèbre
extérieure sur $H^{1}$, et l'annulation en degré $> \dim A_{0}$ borne
$\dim H^{1}$}

\marginal{Si on ne veut pas se servir \uncertain{du résultat} de la torsion, on
\uncertain{trouve} que $\mathrm{Pic}^{\tau}_{X/S}$ \ill{} est simple sur $S$
\ill{} $\mathrm{Pic}^{0}_{A/S}$, \struck{\ill{}} \ill{} voit donc \ill{} qu'il
est un \ill{} surjectif, et de plus \emph{plat}, donc \ill{}, soit le \ill{}
$\mathrm{Pic}^{0}_{A/S}$ est aussi \ill{} dans \ill{} $\mathrm{Pic}_{A/S}$,
c'est donc \ill{} un schéma abélien \ill{} sur $S$ par le
\uncertain{changement} \ill{}}

\note{cette note court en diagonale, à environ $45^{\circ}$, sur toute la marge
gauche du feuillet, sur une vingtaine de lignes. Elle propose une seconde route
vers le théorème, qui éviterait le résultat de torsion invoqué dans le corps du
texte ; c'est la seule chose que l'on puisse en affirmer. Les mots donnés
au-dessus sont ceux que le feuillet soutient ; l'écriture est ici plus rapide
que partout ailleurs dans le dossier et l'essentiel de la liaison est perdu}

\page{3}

et comme $\dim A_{0} \leqslant \dim A'_{0}$ \ill{} un \ill{} général, \ill{} on
a égalité partout ; on trouve en particulier que $A'_{0}$ est \emph{simple}, et
aussi
\[
  \dim H^{1}(A_{0}, \mathcal{O}_{A_{0}}) = n ,
\]
d'où la structure totale de $H^{*}(A_{0}, \mathcal{O}_{A_{0}})$.

Comme on vient de voir que $A'_{0}$ est simple, on voit que pour
$\forall\, x' \in A'_{0}$, $\exists\, x \in A_{0}$ au dessus et un
\uncertain{homomorphisme} \uncertain{vicinaire} $\mathcal{O}_{A'_{0}, x'}
\longrightarrow \mathcal{O}_{A_{0}, x}$ \emph{injectif}.

\note{le qualificatif de l'homomorphisme n'est pas sûr ; ce que l'énoncé demande
est un morphisme d'anneaux locaux injectif}

Or on a le

\textbf{Lemme}\quad $X$ et $Y$ de type fini sur $S$ loc.\ noeth.,
$\varphi : X \to Y$ un $S$-morphisme, $s \in S$,
\struck{$y \in Y$ au dessus} au dessus de $y$, on suppose :

\note{interlinéaire au-dessus de la ligne : \guillemotleft{} (soient $F$, $G$
quasi-coh.\ sur $X$, $Y$ ; $G \to F$ un $\varphi$-morphisme) \guillemotright}

\begin{itemize}
  \item[a)] \struck{$G$} \uncertain{$F$} est plat sur $S_{y}$
    \note{la lettre est réécrite par-dessus une autre} ;
  \item[b/] $G$ coh. ;
  \item[c)] Soit $\varphi_{s} : X_{s} \to Y_{s}$ induit par $\varphi$. Alors
    tout élément de $G_{s}$ dont l'image inverse en les
    $x \in \varphi_{s}^{-1}(y)$ s'annule \ill{}
    \struck{$G_{s} \longrightarrow \varphi_{s*}(F_{s})$ est injectif \ill{} $y$.}
\end{itemize}

\textbf{Alors} $G$ est plat sur $S$ en $y$.

\page{4}

\textbf{Cor.\ Main 1}\quad \struck{$R^{i}f_{*}(\mathcal{O}_{A})$ \ill{}} Soit
\struck{coh.} $A$ plat propre \ill{}, $f$ le
\struck{$\varphi$ morphisme de projection} morphisme de projection
$A \to S$. \struck{L'application} Le morphisme d'Algèbres
\[
  R^{*}f_{*}(\mathcal{O}_{A}) \;\simeq\;
  \Lambda^{*}\!\left(R^{1}f_{*}(\mathcal{O}_{A})\right)
\]
est un isomorphisme, et $R^{1}f_{*}(\mathcal{O}_{A})$ est loc.\ libre sur $S$,
de rang égal à la dimension relative de $A$ sur $S$.

\textbf{Dém.}\quad On sait déjà que \ill{} quand $S$ \ill{} le spectre d'un
corps (cf.\ ci-dessus). D'autre part, il est connu que \ill{} la section unité
$\varepsilon$ implique que
\[
  R^{1}f_{*}(\mathcal{O}_{A}) \;=\;
  \mathrm{Hom}\!\left(\struck{\varepsilon^{*}(\Omega^{1}_{A/S})}\;
  \Omega^{1}_{\varepsilon},\; \mathcal{O}_{S}\right)
\]
\ill{} loc.\ libre sur $S$, et de façon plus précise que $\mathcal{O}_{A}$ est
coh.\ plat pour $f$ en dim.\ 1.

\note{\emph{coh.\ plat} est son abréviation pour \emph{cohomologiquement plat},
qu'il n'écrit nulle part en toutes lettres dans le dossier}

Considérons le \struck{\ill{}} \add{diagramme} canonique :

\begin{tikzcd}[column sep=small, row sep=small, nodes={font=\scriptsize}]
  \Lambda^{*}\!\left(R^{1}f_{*}(\mathcal{O}_{A})\right)_{y}
    \arrow[r, "\alpha"] \arrow[d, "\text{surj}"]
  & R^{*}f_{*}(\mathcal{O}_{A})_{y} \arrow[d] \\
  \Lambda^{*}H^{1}(X_{y}, \mathcal{O}_{X_{y}})
    \arrow[r, "\simeq"]
  & H^{*}(X_{y}, \mathcal{O}_{X_{y}})
\end{tikzcd}

\note{la flèche du bas porte, écrit au-dessus du signe $\simeq$,
\guillemotleft{} Hopf \guillemotright{} ; c'est le théorème de structure des
algèbres de Hopf graduées qui la donne. Il écrit $X_{y}$ pour la fibre de $A$}

\ill{} est \ill{} que
\[
  R^{*}f_{*}(\mathcal{O}_{A})_{y} \longrightarrow H^{*}(X_{y}, \mathcal{O}_{X_{y}})
\]
est surjectif, d'où le fait que $\mathcal{O}_{A}$ est coh.\ plat pour $f$.
Comme $\alpha$ est un hom.\ de Modules \ill{} gradués \ill{} isom.\ \ill{} la
formule \ill{}

\page{5}

réduites, c'est un isomorphisme, cqfd.

\textbf{Remarque}\quad J'ignore si le corollaire vaut sans supposer $A$
projectif sur $S$, et comme j'ignore si $\mathrm{Pic}_{A/S}$ existe dans ce cas,
et s'il \struck{est} \ill{} sur $S$ \ill{}.

\note{c'est la seule ligne du dossier où il dit ce qu'il ne sait pas ; la fin de
la phrase, sur deux mots, ne se laisse pas lire}

\textbf{Corollaire 2}\quad Si $L$ est \uncertain{très ample} rel.\ à $S$, alors
(et non seulement ample) $\varphi_{L} : A \longrightarrow A'$ est un morphisme
surjectif, \struck{\ill{}} et \add{de plus} fidèlement plat ; son noyau est un
gpe fini et plat sur $S$.

\textbf{Corollaire 3}\quad \uncertain{$\mathrm{Pic}^{\tau}_{A/S}$} est ouvert et
\emph{fermé} [l'ouvert est toujours le cas ; le fermé est
\uncertain{spécial} aux schémas abéliens].

\note{l'exposant de $\mathrm{Pic}$ est tracé ici comme une croix, alors que le
$\tau$ des pages 2 et 6 est net}

En effet, si $\varphi_{L}$ \ill{} induit un hom.\ surjectif \ill{} sur $S$, il
induit un hom.\ surjectif \ill{} \ill{} pts sur $\bar{S}$.

\textbf{Corollaire 4}\quad Existence de $\text{Nér-Sév}_{A/S}$.

\page{6}

Nous \ill{} étudier de plus près le morphisme
\[
  (*) \qquad \mathrm{Pic}_{A/S} \longrightarrow \mathrm{Hom}_{S\text{-gr}}(A, A')
\]
i.e.
\[
  (**) \qquad \mathrm{Pic}_{A/S} \times_{S} A \longrightarrow A'
\]
défini par
\[
  (\xi, t) \;\rightsquigarrow\; t_{*}(\xi) - \xi .
\]

\note{à côté de $\mathrm{Hom}_{S\text{-gr}}(A,A')$ il porte, sur deux lignes
reliées par une accolade, \guillemotleft{} hom.\ ponctuel \guillemotright{} et
\guillemotleft{} hom.\ de groupes \guillemotright, avec une flèche de la seconde
vers la première}

\textbf{Théorème}\quad L'hom.\ $(*)$ est un hom.\ de $S$-groupes \uncertain{?}
(i.e.\ $(**)$ est un \uncertain{accouplement}), et son noyau est exactement
\ill{} $A'$ ; \struck{\ill{}} \ill{} section de $\mathrm{Pic}_{A/S}$ \ill{}
\[
  \pi^{*}(\xi) \;=\; \mathrm{pr}_{1}^{*}(\xi) + \mathrm{pr}_{2}^{*}(\xi)
\]

\note{à droite du cadre : \guillemotleft{} $\pi : A \times_{S} A \to A$ \ill{}
la loi de composition \ill{} i.e.\ $\sigma$ \guillemotright. Le noyau est donc
formé des classes \emph{primitives} pour la loi de groupe, ce qui est la
description usuelle de $\mathrm{Pic}^{0}$}

\marginal{$\text{Nér-Sév}_{A/S}
= \mathrm{Pic}_{A/S}/\mathrm{Pic}^{0}_{A/S}$ \ill{} définition}

\note{note de sa main, entourée d'un trait, au bas du feuillet à gauche}

\textbf{Corollaire 6}\quad \ill{}
\[
  \text{Nér-Sév}_{A/S} \longrightarrow
  \mathrm{Hom}_{S\text{-gr}}(A, A')
\]
[compatible avec les structures de groupe.]

\textbf{Démonstration} \ill{} du théorème. L'addition \ill{} est \ill{} au fait
que \ill{} $A \times_{S} A \to A'$ \ill{} par $s$ \ill{}. L'additivité \ill{}
est \ill{} triviale. \ill{} par $(**)$.

\page{7}

\ill{} induit \ill{} sur les \ill{} $A' \times_{S} e$ et $e \times_{S} A$ qui
sont triviaux. Donc il est trivial \ill{} $\xi$ \ill{} un \ill{} sur
$A \times A'$, qui \ill{} un schéma abélien. Donc le \ill{} ramené à prouver
\ill{}

\marginal{ici on utilise que \struck{\ill{}} \ill{} $A'$}

\ill{} le noyau de
\[
  \text{Nér-Sév}_{A/S} \longrightarrow
  \mathrm{Hom}_{S\text{-gr}}(A, A')
\]
\ill{} ensemblistement \ill{}

\textbf{Lemme}\quad $A$ un \ill{} sur \ill{} $S$ \ill{} le spectre d'un \ill{} :
\ill{} pour \ill{} trivialisation. \ill{} $k$ \uncertain{alg.\ clos} \ill{}
$A$, \ill{} Alors $\xi$ est \ill{} $\simeq 0$. Cf.\ \ill{} et sera prouvé plus
bas \ldots

\page{8}

\marginal{L'énoncé \uncertain{intrinsèque} \ill{}}

On a \ill{}, \ill{}
\[
  \varphi_{\xi} = 0 \quad \text{ssi} \quad
  \pi^{*}(\xi) - \mathrm{pr}_{1}^{*}(\xi) - \mathrm{pr}_{2}^{*}(\xi) = 0 .
\]
En effet, considérons le diagramme

\begin{tikzcd}[column sep=small, row sep=small, nodes={font=\scriptsize}]
  A & A \times A \arrow[l, "\mathrm{pr}_{2}"'] \arrow[r, "\pi"]
    \arrow[d, "\mathrm{pr}_{1}"] & A \\
  & A &
\end{tikzcd}

Dire que $\varphi_{\xi} = 0$ signifie aussi \ill{} bien \ill{} que si
$S_{1} = A \to S$, $A_{1} = A \times_{S} S_{1}$, et considérant la \ill{}
section $t_{1}$ de $A_{1}$ sur $S_{1}$,
\[
  t_{1}^{*}(\xi_{1}) = \xi_{1} \qquad (\text{ici } \xi_{1} = \xi \times_{S} S_{1}) .
\]

\struck{\ill{} $t_{1}^{*}(\xi_{1}) = \xi_{1}$ \ill{} par la section unité.}

\ill{} $\xi_{1}$ par $\mathrm{pr}_{2}^{*}(L)$, on a
\[
  t_{1}^{*}\!\left(\mathrm{pr}_{2}^{*}(L)\right) =
  \left(\mathrm{pr}_{2} \circ \mathrm{transl}(t_{1})\right)^{*}(L)
\]
\ill{} est \ill{} \ill{} et qui \ill{} \ill{} définit \ill{} \ill{}. Comme
\ill{} pour \ill{} \ill{}, \ill{} la section unité \ill{}
\[
  i_{1}^{*}(\cdots) = (\pi i_{1})^{*}(L) - (\mathrm{pr}_{1} i_{1})^{*}(L)
  - (\mathrm{pr}_{2} i_{1})^{*}(L) \;=\; 0
\]

\note{sous $\pi i_{1}$ et sous $\mathrm{pr}_{1} i_{1}$ il écrit
\guillemotleft{} id \guillemotright, et sous $\mathrm{pr}_{2} i_{1}$
\guillemotleft{} $\varepsilon_{A/S}$ \guillemotright{} ; les trois accolades
sont ce qui annule le membre de droite}

\page{9}

cette rigidité pour la \ill{} \ill{} quelque \ill{} \uncertain{bien}.

\note{le feuillet ne porte que ces deux lignes, qui achèvent la démonstration
commencée page 8 ; le reste de la page est blanc}

\end{document}
